\begingroup
\ifdefined\TemplateSubsection\let\subsection\TemplateSubsection\fi
\ifdefined\TemplateSubsubsection\let\subsubsection\TemplateSubsubsection\fi
\section{问题三的模型建立与求解}\label{sec:problem3}

\subsection{建模思路与信息边界}
问题三在问题二的日前计划基础上，增加了光伏预报更新与购电合同调整机制。本文将31天历史场景与滚动时域优化结合：0:00利用最新光伏预报制定全天计划；6:00、12:00等更新时刻重新构造剩余时段的净负荷场景，在已执行储能状态和已确认合同的基础上调整未来策略。实际净负荷超出计划净供能的部分，仍由五倍电价紧急购电补足。

将一天按记录顺序划分为144个十分钟时段，$\Delta t=1/6\,\mathrm h$；内部索引$t=1,\ldots,144$，预报时刻$r\in\{0,6,12,18\}$对应已完成时段数$\tau_r=6r$。设$\mathcal T_r=\{\tau_r+1,\ldots,144\}$。每次优化只使用前31天的历史实际数据、当次已发布预报，以及发布前最后一个已完成时段的观测值。目标日未来实际负载和光伏仅用于事后评价，不进入当次决策。

本文延续问题二每日首末储电量为6000 kWh的周期边界，只使用预报中截至当日24:00的部分。日内更新时，初始储电量必须承接已经执行的充放电结果，不能在6:00或12:00重新设为6000 kWh。该边界有利于比较不同预报组合，但限制了跨日储能和次日预报信息的利用。

\subsection{光伏预报插值及净负荷场景}
附件3按每天四个发布时间给出未来24小时整点预报。设$\widehat P_{d,r,h}$为第$d$日$r$时发布的第$h$小时光伏功率预报，令$y_0$为当前已观测光伏功率，$y_h=\widehat P_{d,r,h}$。0:00锚点取零，其他时刻使用发布前最后一条实际记录。对相邻整点间的$x\in[h,h+1]$，线性插值为
\begin{equation}
\widetilde P_{d,r}(x)=(h+1-x)y_h+(x-h)y_{h+1}.
\end{equation}
正文采用端点采样$\widehat P^{10}_{d,r,t}=\widetilde P_{d,r}((t-\tau_r)\Delta t)$，与附件实际功率从0:10至24:00的记录位置配对。为检验实现差异，另与区间中点采样$x=(t-\tau_r-1/2)\Delta t$的已有计算结果比较。两种离散方式对应不同的功率代表值假设，比较时保持各自的预测误差和优化结果完整配套。

对历史日期$s\in\mathcal S_d=\{d-31,\ldots,d-1\}$，按相同发布时间和插值方法计算光伏预测误差，并与同一天负载配对：
\begin{align}
e_{s,r,t}&=P^{\mathrm{PV}}_{s,t}-\widehat P^{10}_{s,r,t},\\
P^{\mathrm{PV},(s)}_{d,r,t}&=\max\{0,\widehat P^{10}_{d,r,t}+e_{s,r,t}\},\\
N^{(s)}_{d,r,t}&=\bigl(P^{\mathrm L}_{s,t}-P^{\mathrm{PV},(s)}_{d,r,t}\bigr)\Delta t,
\qquad \pi_s=1/31.\label{eq:p3-scenario}
\end{align}
式\eqref{eq:p3-scenario}以目标日新预报为中心，叠加近期历史误差，避免将预报视为确定值；非负截断避免出现负光伏功率。负载不另行拟合预测模型，同日负载与光伏误差的配对保留历史联合变化。其适用前提是近期误差具有一定可迁移性，尚不能保证极端天气或负载突变均被31个场景覆盖。

\subsection{初始计划与逐次调整结算}
记$G^0_t$为0:00初始合同购电量；日内更新前后的合同量分别为$G^{\mathrm{old}}_t$和$G^{\mathrm{new}}_t$。定义增购量$A^r_t$和削减量$B^r_t$，则
\begin{equation}
G^{\mathrm{new}}_t=G^{\mathrm{old}}_t+A^r_t-B^r_t,
\qquad A^r_t\geq0,\quad 0\leq B^r_t\leq G^{\mathrm{old}}_t.
\label{eq:p3-contract}
\end{equation}
本文将题面50\%违约电价解释为取消电量仅退回原价的50\%，被取消部分承担其余50\%损失。每次调整相对于上一版已确认合同结算，其净费用为
\begin{equation}
K_d^r=\sum_{t\in\mathcal T_r}p_t(1.5A^r_t-0.5B^r_t).
\label{eq:p3-cash}
\end{equation}
例如原计划100 kWh削减为80 kWh，初始支付$100p$，获得$10p$退款，最终支出$90p$；若增至120 kWh，则支付$100p+1.5p\times20=130p$。因此调整净费用可以为负，不能将退款误记为负违约损失，也不能在保留全部初始费用的同时再次对被取消电量重复收费。此处采用交付时段对应的分时电价$p_t$；逐次结算和退款解释均为本文统一的交易口径。

相同时段若同时增购和削减相同电量，会增加$p_t$倍该电量的费用而不改变合同。因此，正电价下同时增减没有经济优势。多次调整的路径仍影响结算，不能只用最终合同与初始合同之差代替每次交易。

\subsection{随机线性规划及滚动执行}
设$C_t,D_t,E_t$分别为充电输入电量、放电输出电量和时段末储电量，$U^{(s)}_t$为场景紧急购电量。0:00的目标为
\begin{equation}
\min\ \sum_{t=1}^{144}p_tG^0_t
+\frac1{31}\sum_{s\in\mathcal S_d}\sum_{t=1}^{144}5p_tU^{(s)}_t
+\varepsilon\sum_{t=1}^{144}(C_t+D_t),\quad\varepsilon=10^{-7}.
\label{eq:p3-init}
\end{equation}
在$r>0$时，已支付合同费用及过去费用为沉没成本，剩余时域的优化目标改为
\begin{equation}
\min\ \sum_{t\in\mathcal T_r}p_t(1.5A^r_t-0.5B^r_t)
+\frac1{31}\sum_s\sum_{t\in\mathcal T_r}5p_tU^{(s)}_t
+\varepsilon\sum_{t\in\mathcal T_r}(C_t+D_t).
\label{eq:p3-roll}
\end{equation}
各阶段满足共同的设备与场景约束：
\begin{align}
U^{(s)}_t&\geq N^{(s)}_{d,r,t}-G_t-D_t+C_t,\qquad U^{(s)}_t\geq0,\\
E_t&=E_{t-1}+0.9C_t-D_t/0.9,\\
1200&\leq E_t\leq10800,\qquad 0\leq C_t,D_t\leq5000\Delta t,\\
G_t&\geq0,\qquad E_{144}=6000,\\
E_{\tau_r}&=E^{\mathrm{exec}}_{\tau_r}.\label{eq:p3-boundary}
\end{align}
0:00取$E^{\mathrm{exec}}_0=6000$，调整阶段另加入式\eqref{eq:p3-contract}。上述不等式等价于在场景电量平衡中引入非负供能富余量后将其消去。每个阶段的合同与充放电变量在31个场景间共享，紧急购电变量随场景变化。目标中的小惩罚仅用于减少不必要的吞吐；0:00代码另对购电量加入$10^{-9}$系数增量，报告费用均剔除这些数值辅助项。

剩余时段数为$H$时，初始模型含$35H+1$个变量，调整模型含$36H+1$个变量，均有$H+2$条等式及$31H$条场景不等式。通过SciPy调用HiGHS求解，按时域长度缓存稀疏矩阵。对所选预报集合$\mathcal R=\{0,6,12\}$，实际执行遵循
\begin{equation}
\underbrace{\text{0点求解全天}}_{\text{执行0至6点}}
\longrightarrow
\underbrace{\text{6点重算余下18小时}}_{\text{执行6至12点}}
\longrightarrow
\underbrace{\text{12点重算余下12小时}}_{\text{执行12至24点}}.
\end{equation}
每次只冻结执行到下一个更新时刻的结果，未执行部分可以调整，式\eqref{eq:p3-boundary}保证储能状态连续。该方法是滚动求解一系列两阶段随机规划，并未建立包含未来预报分支的完整多阶段场景树，因此单次LP最优不等于全年实际费用全局最优。

\subsection{实际费用与实现差异比较}
将各阶段已执行片段拼接为$G^{\mathrm{exec}}_{d,t},C^{\mathrm{exec}}_{d,t},D^{\mathrm{exec}}_{d,t}$，使用实际净负荷$N_{d,t}$计算
\begin{align}
U^{\mathrm{act}}_{d,t}&=[N_{d,t}-G^{\mathrm{exec}}_{d,t}-D^{\mathrm{exec}}_{d,t}+C^{\mathrm{exec}}_{d,t}]^+,\\
F_d&=\sum_tp_tG^0_{d,t}+\sum_{r\in\mathcal R\setminus\{0\}}K_d^r
+\sum_t5p_tU^{\mathrm{act}}_{d,t}.\label{eq:p3-total}
\end{align}
实际外网购电量为$\sum_t(G^{\mathrm{exec}}_{d,t}+U^{\mathrm{act}}_{d,t})$，与0:00合同量及最终合同量均有区别。富余供能不获得售电收入，其中可能包含未消纳合同电量，不能全部解释为弃光。

对已有端点与中点两种实现进行比较。两者的31天场景、储能参数、逐次退款规则及日末边界相同，主要区别在预报采样位置、富余变量是否显式保留，以及求解器退化处理。中点实现采用显式场景平衡与更严格的对偶容差，端点实现消去富余变量并在初始购电项加入微小扰动。结果见表\ref{tab:p3-versions}，所有费用均覆盖2025年2月1日至12月31日334天。
\begin{table}[htbp]\centering
\caption{两种插值实现的同策略结果比较}\label{tab:p3-versions}
\small\begin{tabular}{lrrr}\toprule
实现与预报组合 & 总费用/元 & 紧急购电量/kWh & 调整净费用/元\\\midrule
端点，$0+6+12$（正文）&17013832.1962&281014.6588&6374.8705\\
中点，$0+6+12$ &17020306.6465&282085.6383&9377.4177\\
中点，$0+6+12+18$ &17020342.9202&282112.5626&9278.6275\\\bottomrule
\end{tabular}\end{table}
端点方案在相同三次预报组合下少支出6474.4504元，差异为中点方案总费用的0.0380\%。本文据其记录配对方式、可复现性及本次回测费用选用端点方案；中点方案作为离散口径对照。该差异同时受场景采样和数值实现影响，不能解释为某种插值方法的普遍优势。已有四次预报主结果不直接作为三次预报图表使用；正文全部指定日期结果、年度统计与执行图均来自同一端点三次预报策略。

\subsection{指定日期购电与储能结果}
表\ref{tab:p3-plan}至表\ref{tab:p3-emergency}按题面表1至表3结构分别列出初始及调整后合同购电、最终充放电和紧急购电情况。购电结果按模板时间标签逐列对应，例如10:00记录填入10:00--10:10；四小时充放电量按连续24条记录求和。与问题二一致，模板标签相对内部时段起止存在十分钟边界差异，本文保留原结果映射，计算与滚动执行均以内部索引为准，未据标签移动数值。

四个日期的实际总费用依次为48029.3636、46014.0495、47603.5756和65428.9947元。3月20日、9月23日及12月21日的调整净费用分别为$-1255.9632$、$-562.5049$和$-58.4863$元，表示削减合同的退款超过增购支出；6月21日调整净支出69.4475元，全天无紧急购电。3月20日紧急购电1237.0904 kWh，高于问题二对应日的207.4009 kWh，说明全期费用改善不意味着每一天的补购量都下降。从原始附件重算这四天，合同曲线与已有解一致，费用差异不超过$1.456\times10^{-11}$元。

\input{sections/problem3_tables}

图\ref{fig:p3-dispatch}以12月21日展示合同、紧急购电、充放电及储电量的对应关系。6:00和12:00更新后，仅剩余时段合同与储能计划变化，储电量沿原状态连续演化。该日紧急购电1398.7836 kWh，费用4448.9943元，表明更新预报仍不能消除所有实际偏差。图中富余量和放电量向下绘制仅为区分方向；真实变量均非负。其余三个指定日期的同口径执行图保存在支撑材料中。
\begin{figure}[htbp]\centering
\includegraphics[width=0.96\textwidth]{figures/problem3/dispatch_2025-12-21.pdf}
\caption{12月21日三次预报策略执行结果；竖虚线为6:00与12:00更新点}
\label{fig:p3-dispatch}
\end{figure}

\subsection{其他预报时刻的引入价值}
\subsubsection{基于实际费用的消融比较}
固定历史窗口、日末边界和结算规则，分别重算六种预报组合，结果如表\ref{tab:p3-ablation}。仅使用0:00预报作为问题三内部基准；它与不使用附件3的问题二基准不同。
\begin{table}[htbp]\centering
\caption{不同预报组合的334天回测结果（端点插值）}\label{tab:p3-ablation}
\small\begin{tabular}{lrrr}\toprule
预报组合 & 总费用/元 & 紧急购电量/kWh & 相对仅0点节省/元\\\midrule
$0$&17361854.7485&379528.7227&0.0000\\
$0+6$&17096655.4908&299802.2818&265199.2576\\
$0+6+12$&\textbf{17013832.1962}&\textbf{281014.6588}&\textbf{348022.5523}\\
$0+6+12+18$&17013912.4490&281047.6368&347942.2995\\
$0+12+18$&17109751.5453&321831.1428&252103.2032\\
$0+6+18$&17096348.1608&299814.4209&265506.5877\\\bottomrule
\end{tabular}\end{table}
沿$\{0\}\to\{0,6\}\to\{0,6,12\}\to\{0,6,12,18\}$依次引入预报，6:00节省265199.2576元，12:00进一步节省82823.2947元，18:00则增加80.2528元。三次预报相比仅0:00降低费用2.0045\%。从四次预报组合分别删除6:00、12:00、18:00，费用变化为增加95839.0963元、增加82435.7118元和减少80.2528元。因此，逐次加入与逐项删除两种比较均支持保留6:00、12:00更新。

图\ref{fig:p3-value}采用相对基准的节省额展示六种组合，并将逐次增加预报的边际节省单独绘制，使18:00的微小差异不被千万元级总费用掩盖。中点实现的6:00和12:00同样具有正边际价值，而加入18:00增加36.2737元，说明三次预报的选择在两种实现下相同。
\begin{figure}[htbp]\centering
\includegraphics[width=\textwidth]{figures/problem3/ablation_value.pdf}
\caption{预报组合的经济性比较；左图分别使用各插值实现自身的0点基准，右图为端点结果}
\label{fig:p3-value}\end{figure}

加入18:00后，端点方案调整净支出减少101.7770元，但紧急购电费用增加182.0298元，净效应为增加80.2528元。该数额相对全期费用极小，结论应表述为“当前回测中未见额外经济收益”，而非18:00预报必然无用。六种组合也未穷举所有可选子集，且预报组合在同一年度回测上选择，尚需跨年度或独立验证期确认泛化效果。

\subsubsection{预报误差与经济价值的区别}
用MAE和RMSE评价光伏功率预报，避免夜间实际功率为零时MAPE失真。各发布时间针对其剩余时段的RMSE依次为383.0269、271.3714、139.7280和42.9234 kW，但相应样本数依次为48096、36072、24048和12024，不能直接把误差下降全部归因为预报更新。

为公平比较，进一步固定18:00至24:00的共同窗口，每个预报均使用12024个样本。对应RMSE为46.2420、43.9872、43.4104和42.9234 kW，12:00至18:00仅改善0.4870 kW。图\ref{fig:p3-accuracy}显示，共同晚间窗口的误差差异远小于各自剩余窗口的差异。即使预报误差略有下降，也需通过优化后紧急购电与合同调整的合计费用判断其价值。
\begin{figure}[htbp]\centering
\includegraphics[width=\textwidth]{figures/problem3/forecast_accuracy.pdf}
\caption{各预报时刻的光伏误差；左图窗口不同，右图统一为18:00至24:00}
\label{fig:p3-accuracy}\end{figure}

\subsection{全期结果与可行性检验}
所选三次预报策略的全期结果见表\ref{tab:p3-summary}，其中“全期”均指附件5要求的2月至12月334天，1月用于初始化31天历史窗口。

\begin{table}[H]\centering
\caption{三次预报策略的334天购电汇总}\label{tab:p3-summary}
\small\setlength{\tabcolsep}{4pt}
\begin{tabular}{lrlr}\toprule 指标&结果&指标&结果\\\midrule
初始合同量/kWh&25629550.9844&初始计划费/元&15936251.7627\\
执行合同量/kWh&24998657.1288&调整净费用/元&6374.8705\\
紧急购电量/kWh&281014.6588&紧急购电费/元&1071205.5630\\
实际总购电量/kWh&25279671.7876&实际总费用/元&17013832.1962\\
紧急电量占比&1.1116\%&日均总费用/元&50939.6174\\
紧急购电天数/天&243&计算天数/天&334\\\bottomrule
\end{tabular}\end{table}

相较问题二，实际总费用减少595959.4904元，降幅3.3843\%；紧急购电量减少196431.3986 kWh，降幅41.1421\%。其中，从问题二到仅使用0:00光伏预报，费用减少247936.9381元；再加入6:00和12:00滚动调整，减少348022.5523元。该分解区分了初始预报信息与日内调整的收益，不能将总收益全部归于滚动控制。

执行合同量比初始合同量净减少630893.8556 kWh，但调整净费用仍为正，反映增购和削减在不同时段发生且结算单价不对称。仍有243天发生紧急购电，占334天的72.7545\%；供能富余总量为5939095.8012 kWh。因此该策略主要降低费用与补购规模，并未消除缺口发生。供能富余包含合同富余，不据此计算弃光率。

对48096个时段逐一核验，供需平衡最大绝对残差为$2.274\times10^{-13}$ kWh，储能状态转移残差为$9.095\times10^{-13}$ kWh，首末储电量偏差不超过$5.457\times10^{-12}$ kWh；合同更新恒等式残差不超过$2.274\times10^{-13}$ kWh，逐日费用重算残差不超过$2.388\times10^{-11}$元。储电量满足1200至10800 kWh边界，单时段充、放电上限均为833.3333 kWh；以$10^{-7}$ kWh为阈值检查，同时充放电的时段数为0。虽未显式加入互斥二进制变量，当前数值解通过了事后互斥检验。

\endgroup
